\begin{table}[htbp]
\centering
\footnotesize
\caption{Problema G (google/gemini-2.5-flash-lite): nomes de classe semânticos versus neutros.}
\label{tab:bloco3_G_nomes_semanticos_vs_neutros_flashlite}
\begin{tabular}{llrrrrrrrr}
\toprule
 &  & \multicolumn{5}{c}{Acur.\ LLM vs real (\%)} & \multicolumn{3}{c}{Acur.\ métrica vs real (\%)} \\ \cmidrule(lr){3-7} \cmidrule(lr){8-10}
Subproblema & Atributos & M/H literal & A/B literal & M/H inv.\ perm. & A/B inv.\ perm. & $\Delta$ & M/H inv.\ perm. & A/B inv.\ perm. & $\Delta$ \\
\midrule
G1 (normalizado) & peso/altura & 35,2 & 18,6 & 64,8 & 81,4 & $-$16,7 & 89,5 & 83,3 & +6,2 \\
G1 (normalizado) & x1/x2 & 57,1 & 16,2 & 57,1 & 83,8 & $-$26,7 & 74,8 & 90,5 & $-$15,7 \\
\midrule
G2 (kg e m) & peso\_kg/altura\_m & 84,3 & 46,7 & 84,3 & 54,3 & +30,0 & 91,4 & 77,1 & +14,3 \\
G2 (kg e m) & x1/x2 & 48,6 & — & 51,4 & — & — & 65,2 & — & — \\
\bottomrule
\multicolumn{10}{p{0.95\linewidth}}{\footnotesize Efeito do nome das classes (M/H = Mulher/Homem versus A/B) para cada par de nomes de atributo, zero-shot, R$^2$, Perceptron. Acurácia (vs rótulo real) do LLM e da métrica, contexto. Com nomes neutros (e em zero-shot) a correspondência entre nome e classe é arbitrária: a acurácia literal pode sair invertida; as colunas ``inv.\ perm.'' usam o valor invariante à permutação max($a$, 1 $-$ $a$), e $\Delta$ = semântico $-$ neutro é calculado sobre elas. ``—'' = colapso em classe única.} \\
\end{tabular}
\end{table}
